%% ma: L11-B23-0002
%% lop: 11
%% chuong: 7
%% bai: 23
%% dang: 11.23.2
%% loai: TN4
%% muc: TH
%% dap_an: "A"
%% nguon: "(NBV)-ĐỀ SỐ 10-MỖI NGÀY 1 ĐỀ THI 2025.pdf (D10, MỖI NGÀY 1 ĐỀ - Đề số 10), Phần I Câu 4"
%% kien_thuc: "Bài 23 (đường thẳng vuông góc với mặt phẳng), Bài 22 (góc giữa hai đường thẳng), tỉ số lượng giác trong tam giác vuông"
%% trang_thai: chinh-thuc
%% ghi_chu: "Đã duyệt 11/10/2026. Gốc không có hình vẽ. Câu xếp ở Bài 23 vì lời giải dùng SA vuông góc với mặt phẳng đáy suy ra SA vuông góc với AD."
\begin{ex}
Cho hình chóp $S.ABCD$ có đáy $ABCD$ là hình chữ nhật. Biết $SA=AB=a$, $AD=a\sqrt3$. Cạnh bên $SA\perp(ABCD)$. Số đo góc giữa hai đường thẳng $SD$ và $BC$ là
\choice{\True $30^\circ$}{$45^\circ$}{$90^\circ$}{$60^\circ$}
\loigiai{Vì $BC\parallel AD$ nên góc giữa $SD$ và $BC$ bằng góc giữa $SD$ và $AD$, tức $\widehat{SDA}$ (tam giác $SAD$ vuông tại $A$ vì $SA\perp(ABCD)$). Ta có $\tan\widehat{SDA}=\dfrac{SA}{AD}=\dfrac{a}{a\sqrt3}=\dfrac1{\sqrt3}$ nên $\widehat{SDA}=30^\circ$.}
\end{ex}
